\documentclass[11pt]{article}

\usepackage[T1]{fontenc}
\usepackage[margin=1in]{geometry}
\usepackage{amsmath, amssymb}
\usepackage{booktabs}
\usepackage{graphicx}
\usepackage{natbib}
\usepackage{hyperref}
\usepackage{microtype}
\usepackage{placeins}
\usepackage{xcolor}

\hypersetup{
  colorlinks = true,
  linkcolor = blue,
  citecolor = blue,
  urlcolor = blue
}

\newcommand{\R}{\mathbb{R}}
\newcommand{\given}{\,\vert\,}
\newcommand{\ind}{\mathbb{I}}
\newcommand{\mcP}{p_{\mathrm{MC}}}
\newcommand{\code}[1]{\texttt{\small #1}}
\newcommand{\maybeincludegraphics}[2][]{%
  \IfFileExists{#2}{%
    \includegraphics[#1]{#2}%
  }{%
    \fbox{\parbox{0.9\linewidth}{Figure file not found: \texttt{#2}}}%
  }%
}
\setlength{\emergencystretch}{3em}

\title{Beyond the Next Step for Latent-State Movement Models: Generative Validation Diagnostics for Hidden Markov Step-Selection Functions}
\author{Aurelien Nicosia}
\date{\today}

\begin{document}

\maketitle

% The introduction and conclusion are intentionally omitted in this first draft.
% The current version starts with the methodological core and empirical material.

\section*{Scope of this draft}

This draft focuses on the methodological core of the article: the HMM-SSF
model formulation, the generative-validation logic, the diagnostic criteria,
and the first simulation and empirical examples. The introduction, broader
discussion, and conclusion are intentionally left for a later version, after
the final simulation design and empirical emphasis have been fixed.

\section{Extension of the SSF generative-validation framework}

This manuscript is conceived as a latent-state extension of the generative
validation framework for SSF and iSSF models developed in \citet{nicosia2026}.
The starting point is the same: a fitted step-selection model should not only
describe local one-step decisions, but should also regenerate plausible
trajectory-scale patterns when iterated forward. In the single-state setting,
the validation target is the distribution of simulated trajectories generated
by the fitted SSF or iSSF kernel. In the HMM-SSF setting, the validation target
has two coupled components: the spatial trajectory and the latent regime
sequence that governs state-specific movement and selection.

The original framework is organized around four trajectory-level pillars:
emergent utilization distributions, mean squared displacement, path sinuosity,
and barrier crossing. These pillars remain valid for HMM-SSFs because a
multi-state model still generates an ordinary movement trajectory. The HMM
extension adds a second diagnostic layer for the latent regime process:
occupancy, residence times, switching rates, transition counts, and
state-conditioned movement summaries. The goal is therefore not to replace the
SSF diagnostics, but to decompose generative mismatch into spatial,
movement-dynamic, and state-process components.

\begin{table}[ht]
\centering
\caption{Relationship between the original SSF generative-validation pillars and the HMM-SSF extension.}
\label{tab:ssf-to-hmm-extension}
\begin{tabular}{p{0.25\linewidth}p{0.31\linewidth}p{0.34\linewidth}}
\toprule
Original SSF pillar & Role retained for HMM-SSFs & Additional latent-state question \\
\midrule
Emergent utilization distributions & Does the full fitted model regenerate the observed spatial footprint? & Are discrepancies driven by the state-specific SSF kernels or by unrealistic state occupancy? \\
Mean squared displacement & Does the model reproduce displacement scaling over time? & Does the latent switching process reproduce the temporal organization of movement regimes? \\
Path sinuosity & Does the model reproduce path-level tortuosity or directional persistence? & Are state-specific movement kernels and state durations jointly producing realistic path structure? \\
Barrier crossing & Does the model reproduce functional permeability of known barriers? & Does barrier response vary by state, and are barrier-crossing opportunities assigned to realistic regimes? \\
\bottomrule
\end{tabular}
\end{table}

This framing also clarifies how the diagnostics should be interpreted. The
Monte Carlo p-values are not generic measures of visual closeness across
candidate models. Each fitted model generates its own reference distribution,
with its own variability and sharpness. A model can therefore have a larger
observed discrepancy but a less extreme p-value if its simulated trajectories
are highly variable, whereas a sharper model can yield a smaller p-value for a
smaller absolute discrepancy. The diagnostic question is model-specific:
is the observed trajectory unusual under the generative distribution implied by
that fitted model?

\section{Model framework}

Let $\{X_t\}_{t=1}^T$ denote a regular movement trajectory, with
$X_t \in \R^2$, and let $\{S_t\}_{t=1}^T$ denote a latent behavioural regime
process taking values in $\{1,\ldots,K\}$. A hidden Markov step-selection
function (HMM-SSF) combines a latent state process with a state-dependent
movement and habitat-selection kernel. The central modelling assumption is
that the animal alternates among a finite number of movement regimes, and that
each regime has its own movement kernel and habitat-selection response.

\paragraph{Latent-state process.}

The latent process is specified through an initial distribution
$\delta=(\delta_1,\ldots,\delta_K)$ and transition probabilities
\begin{equation*}
  \Pr(S_{t+1}=j \given S_t=i, z_t) = \Gamma_{ij}(z_t),
  \qquad i,j \in \{1,\ldots,K\},
\end{equation*}
where $z_t$ denotes optional covariates affecting state persistence or
switching. In the homogeneous case, $\Gamma_{ij}(z_t)$ reduces to a fixed
transition matrix $\Gamma$. In a time-varying transition model, each row of
$\Gamma(z_t)$ is usually parametrized by multinomial-logit contrasts. For
example, using state $i$ as the baseline destination from current state $i$,
one may write
\begin{equation*}
  \log\left\{\frac{\Gamma_{ij}(z_t)}{\Gamma_{ii}(z_t)}\right\}
  =
  \alpha_{ij} + \lambda_{ij}^{\top} z_t,
  \qquad j \ne i.
\end{equation*}
This representation separates behavioural persistence from movement and
habitat selection. The transition model controls when regimes are entered and
left, whereas the state-dependent SSF controls how the animal moves and selects
habitat while occupying a given regime.

\paragraph{State-dependent step-selection kernel.}

Conditional on the next state, the transition from $X_t$ to $X_{t+1}$ is
described by a state-specific step-selection kernel. In its exponential
form, a generic state-specific SSF can be written as
\begin{equation}
  f_s(x \given x_t, x_{t-1}, e_t)
  =
  \frac{
    q_s(x \given x_t, x_{t-1}; \eta_s)
    \exp\{ \beta_s^\top c(x, x_t, x_{t-1}, e_t) \}
  }{
    \int_{\Omega}
    q_s(u \given x_t, x_{t-1}; \eta_s)
    \exp\{ \beta_s^\top c(u, x_t, x_{t-1}, e_t) \}
    \,du
  },
  \label{eq:ssf-kernel}
\end{equation}
where $q_s$ is the movement kernel in state $s$, $\eta_s$ contains its
movement parameters, $e_t$ denotes environmental information available at time
$t$, $c(\cdot)$ is a vector of movement and environmental covariates, and
$\beta_s$ is a state-specific selection coefficient vector. This formulation
follows the general step-selection framework used to connect movement
constraints and resource selection \citep{forester2009,avgar2016} and its
multi-state HMM-SSF extensions \citep{nicosia2017,klappstein2023}.

In applications fitted with used and available steps, the continuous
normalizing integral in Eq.~\eqref{eq:ssf-kernel} is replaced by a discrete
normalization over a candidate set $\mathcal{A}_t$ that contains the observed
endpoint and a finite number of available endpoints. The state-specific
probability assigned to candidate $a \in \mathcal{A}_t$ is then
\begin{equation}
  p_s(a \given \mathcal{A}_t)
  =
  \frac{
    q_s(a \given x_t, x_{t-1}; \eta_s)
    \exp\{ \beta_s^\top c(a, x_t, x_{t-1}, e_t) \}
  }{
    \sum_{u \in \mathcal{A}_t}
    q_s(u \given x_t, x_{t-1}; \eta_s)
    \exp\{ \beta_s^\top c(u, x_t, x_{t-1}, e_t) \}
  }.
  \label{eq:discrete-ssf}
\end{equation}
This discrete form is useful for fitting because it connects the HMM-SSF to
conditional logistic regression and to state-dependent conditional likelihood
contributions.

\paragraph{Joint generative model.}

Conditioning on the first locations required to define step length and turning
angle, the joint density of a latent state path and a movement path factorizes
as
\begin{equation}
  p_{\theta}(s_{1:T}, x_{1:T})
  =
  \delta_{s_1}
  \prod_{t=1}^{T-1}
  \Gamma_{s_t s_{t+1}}(z_t)
  f_{s_{t+1}}(x_{t+1} \given x_t, x_{t-1}, e_t),
  \label{eq:joint-generative}
\end{equation}
up to the conditioning convention for the first step. The parameter vector
$\theta$ contains the initial distribution, transition parameters,
state-specific movement parameters, and state-specific selection coefficients.
Equation~\eqref{eq:joint-generative} is the key theoretical object for this
paper: it defines both the likelihood-based model and the simulation mechanism
used for prediction and generative validation.

\paragraph{Likelihood and state uncertainty.}

Because the state path is latent, the observed-data likelihood sums over all
possible state sequences. Let $E_t(s)$ denote the state-specific contribution
of the observed step at time $t$ under state $s$. In the discrete used-available
form, $E_t(s)$ is given by Eq.~\eqref{eq:discrete-ssf} evaluated at the
observed endpoint. The likelihood can then be evaluated with the standard HMM
forward recursion \citep{rabiner1989,langrock2012}. With
$P_t=\mathrm{diag}\{E_t(1),\ldots,E_t(K)\}$, a schematic likelihood expression
is
\begin{equation*}
  L(\theta)
  =
  \delta^\top P_1
  \Gamma(z_1) P_2
  \Gamma(z_2) \cdots
  \Gamma(z_{T-1}) P_T
  \mathbf{1},
\end{equation*}
where $\mathbf{1}$ is a vector of ones. The exact indexing of $P_t$ depends on
whether the first step is treated as fixed or modelled explicitly, but the
factorization is the same.

This likelihood representation matters for validation because state labels are
not directly observed in empirical telemetry data. The fitted model provides
filtered state probabilities, smoothed state probabilities, and a most likely
Viterbi path. These summaries are useful, but they are not equivalent to
observed regimes. Consequently, any diagnostic that uses state labels must be
interpreted conditionally on a state-assignment rule or averaged over plausible
state paths.

\paragraph{Prediction targets.}

A fitted HMM-SSF induces several distinct predictive objects. The one-step
predictive distribution is a mixture over possible next states:
\begin{equation}
  p_{\theta}(x_{t+1} \given \mathcal{F}_t)
  =
  \sum_{i=1}^{K}\sum_{j=1}^{K}
  \alpha_t(i)\Gamma_{ij}(z_t)
  f_j(x_{t+1} \given x_t,x_{t-1},e_t),
  \label{eq:one-step-predictive}
\end{equation}
where $\mathcal{F}_t$ denotes the information available up to time $t$, and
$\alpha_t(i)=\Pr(S_t=i\given \mathcal{F}_t)$ is the filtered probability of
state $i$. Multi-step prediction is obtained by iterating the same mechanism:
forecast the next state, draw or integrate over the next location, update the
movement history and environmental covariates, and repeat. This iteration is
what links local step-selection rules to emergent predictions such as
long-term space use, displacement scaling, path structure, residence times,
and barrier-crossing frequencies.

The distinction between local and generative prediction is central. A model may
provide reasonable one-step conditional probabilities while producing
unrealistic long-horizon trajectories once Eq.~\eqref{eq:joint-generative} is
iterated. The diagnostics below are designed to evaluate this longer-horizon
predictive distribution rather than only the local likelihood contributions.

\section{Generative validation}

The validation workflow uses the same Monte Carlo rank logic as the SSF
framework in \citet{nicosia2026}. It asks whether the observed trajectory, or a
decoded state sequence derived from it, is unusual relative to replicate
outputs generated by the fitted model. Let $y$ denote the observed trajectory
and let $Y^{(1)},\ldots,Y^{(M)}$ denote $M$ simulated trajectories generated
under a fitted HMM-SSF. For a diagnostic discrepancy $D$, the observed
discrepancy is compared with a simulation reference distribution:
\begin{align*}
  T_{\mathrm{obs}}
  &=
  \frac{1}{M}\sum_{m=1}^{M} D(y, Y^{(m)}), \\
  T_m
  &=
  \frac{1}{M-1}\sum_{\ell \ne m} D(Y^{(m)}, Y^{(\ell)}),
  \qquad m=1,\ldots,M .
\end{align*}
The Monte Carlo p-value is
\begin{equation*}
  \mcP =
  \frac{1 + \sum_{m=1}^{M} \ind\{T_m \ge T_{\mathrm{obs}}\}}
       {M+1}.
\end{equation*}
For diagnostics where larger discrepancies indicate poorer agreement, small
values of $\mcP$ indicate that the observed trajectory or decoded state path is
more discrepant than expected under the fitted model. The finite simulation
resolution is $(M+1)^{-1}$, which is reported with each diagnostic.

This rank-based construction is related to simulation-based model criticism and
realized-discrepancy checks \citep{gelman1996} and to Monte Carlo significance
testing \citep{besag1989}. The interpretation is conditional on the fitted
model, the simulation strategy, and the chosen discrepancy. It is not a direct
ranking of candidate models by visual distance, because each candidate model
induces its own simulation reference distribution.

Compatibility alone is not sufficient for an informative diagnostic. A fitted
model can pass because the simulated reference distribution is diffuse. The
workflow therefore reports a sharpness score, computed from the concentration
of simulated discrepancies. In the workflow used here, sharpness is a
secondary label: it does not replace the Monte Carlo p-value, but it helps
distinguish a concentrated compatibility check from one with low diagnostic
resolution.

\section{Diagnostic discrepancies for HMM-SSFs}

The single-state generative-validation framework organizes model checking
around biologically interpretable trajectory-level summaries: emergent space
use, displacement dynamics, path structure, and barrier interactions. These
criteria remain necessary for HMM-SSFs because the fitted model still generates
ordinary spatial trajectories. However, latent-state models introduce an
additional process that must also be checked: the temporal organization of
behavioural regimes. A multi-state model may reproduce the spatial footprint
of the observed trajectory while assigning unrealistic time budgets to states,
generating implausible residence times, or producing state-specific movement
patterns that do not match the decoded regimes.

We therefore separate the diagnostics into two complementary groups. The first
group applies the original trajectory-level checks to the full simulated path,
without conditioning on state labels. The second group targets the latent
regime process and asks whether the fitted model regenerates plausible
occupancy, persistence, switching structure, and state-conditioned movement.
These diagnostics are not combined into a single scalar goodness-of-fit
measure. As in the one-state framework, they are intended to form a diagnostic
dashboard in which each criterion probes a distinct way in which the fitted
model can fail to generate realistic movement.

\begin{table}[ht]
\centering
\caption{Overview of generative diagnostics for HMM-SSFs. The first four rows extend the trajectory-level pillars of the single-state SSF framework; the remaining rows target the latent state process and state-conditioned movement.}
\label{tab:hmm-diagnostics-overview}
\renewcommand{\arraystretch}{1.18}
\setlength{\tabcolsep}{3pt}
\footnotesize
\begin{tabular}{p{0.17\linewidth}p{0.28\linewidth}p{0.24\linewidth}p{0.18\linewidth}}
\toprule
Diagnostic family & Target pattern & Summary statistic & Typical failure modes \\
\midrule
Emergent utilization distributions &
Long-term spatial configuration of space use. &
Spatial discrepancy between observed and simulated utilization distributions. &
Spatial leakage, over-diffusion, or failure to reproduce concentration. \\

Mean squared displacement &
Temporal displacement scaling across lags. &
Integrated discrepancy between observed and simulated MSD curves. &
Over-diffusion, under-diffusion, or missing multi-state temporal structure. \\

Path sinuosity &
Path-level tortuosity and directional persistence. &
Absolute discrepancy in straightness index. &
Paths too straight, too tortuous, or behaviourally unrealistic. \\

Barrier crossing &
Functional connectivity across known linear features. &
Observed barrier crossings relative to simulated crossings. &
Overestimated permeability or state-specific barrier response mismatch. \\

State occupancy &
Time budget allocated to each latent regime. &
Per-state occupancy and global total-variation discrepancy. &
Dominant states, missing states, or unrealistic regime allocation. \\

Residence time &
Persistence and dwell-time shape within regimes. &
Residence-time distribution discrepancy and geometric dwell-time check. &
Non-geometric residence times, memory, or inadequate Markov persistence. \\

Switching and transitions &
Frequency and direction of transitions among regimes. &
Switching rate and normalized transition-count discrepancy. &
Incorrect switching frequency or misspecified time-varying transition model. \\

State-conditioned movement &
Movement phenotype within each decoded or simulated regime. &
State-specific step-length, turning-angle, and MSD discrepancies. &
Misspecified state-specific movement kernels or poorly separated regimes. \\
\bottomrule
\end{tabular}
\end{table}

\subsection{Trajectory-level diagnostics}

Trajectory-level diagnostics do not require a latent state path. They evaluate
emergent movement patterns produced by the full fitted model, and correspond
to the original SSF generative-validation pillars. Let
$X_{1:T}^{(0)}$ denote the observed trajectory and
$X_{1:T}^{(m)}$, $m=1,\ldots,M$, denote simulated trajectories generated under
a given HMM-SSF simulation strategy. These diagnostics intentionally ignore the
state labels after simulation: they ask whether the fitted model, after
integrating over its latent regime process, produces realistic spatial paths.
This distinction is important because a model can have interpretable latent
states while still producing unrealistic long-horizon trajectories, or can
generate plausible spatial paths while allocating movement to unrealistic
regimes.

\paragraph{Utilization-distribution discrepancy.}
This diagnostic targets emergent space use, as in the original SSF framework.
For each trajectory, we construct an empirical utilization distribution (UD)
from the visited locations. Let $\mu_0$ denote the empirical UD of the observed
trajectory and let $\mu_m$ denote the empirical UD of simulated trajectory
$m$. In the full spatial version, mismatch between UDs can be measured by the
1-Wasserstein distance, which quantifies the minimal amount of spatial
``work'' required to transport probability mass from one distribution to
another \citep{villani2009}. In the present HMM-SSF analyses, we use a
computationally light approximation based on the sum of one-dimensional
Wasserstein distances between marginal coordinate distributions:
\begin{equation*}
  D_{\mathrm{UD}}(a,b)
  =
  W_1(a_x,b_x) + W_1(a_y,b_y),
\end{equation*}
where $a_x$ and $b_x$ are the empirical distributions of the $x$ coordinates,
and $a_y$ and $b_y$ are the corresponding empirical distributions of the $y$
coordinates. The observed discrepancy is
\begin{equation*}
  T_{\mathrm{obs}}^{\mathrm{UD}}
  =
  \frac{1}{M}\sum_{m=1}^{M} D_{\mathrm{UD}}(\mu_0,\mu_m),
\end{equation*}
and the simulated reference discrepancies are computed by leave-one-out
comparison among simulated trajectories:
\begin{equation*}
  T_m^{\mathrm{UD}}
  =
  \frac{1}{M-1}
  \sum_{\ell \ne m} D_{\mathrm{UD}}(\mu_m,\mu_\ell).
\end{equation*}
Under a well-specified generative model, the observed UD should not be more
spatially isolated from the simulated UDs than simulated UDs are from one
another. A small Monte Carlo p-value therefore indicates that the fitted model
does not reproduce the long-run spatial footprint of the observed trajectory.
In an HMM-SSF, this failure can arise from the state-specific SSF kernels, from
unrealistic state occupancy or residence near particular habitats, or from an
interaction between movement and switching dynamics.

\paragraph{Mean squared displacement discrepancy.}
Mean squared displacement (MSD) targets the temporal scaling of movement
\citep{turchin1998}, and is closely related to net squared displacement
summaries used to characterize movement phenotypes \citep{bastilleRousseau2016}.
For trajectory $r \in \{0,1,\ldots,M\}$ and lag $k$, define
\begin{equation*}
  \mathrm{MSD}_r(k)
  =
  \frac{1}{T-k}\sum_{t=1}^{T-k}
  \|X_{t+k}^{(r)}-X_t^{(r)}\|^2 .
\end{equation*}
Let $\mathcal{K}$ be the set of lags used for the diagnostic, and let
$w_k \ge 0$ be optional weights used to emphasize particular temporal scales
(with $w_k=1$ in the unweighted case). The model-expected MSD curve is
estimated from the simulated ensemble,
\[
  \overline{\mathrm{MSD}}_{\mathrm{sim}}(k)
  =
  \frac{1}{M}\sum_{m=1}^{M}\mathrm{MSD}_m(k).
\]
The observed MSD discrepancy is then
\begin{equation}
  T_{\mathrm{obs}}^{\mathrm{MSD}}
  =
  \sum_{k \in \mathcal{K}}
  w_k
  \left[
    \mathrm{MSD}_0(k)
    -
    \overline{\mathrm{MSD}}_{\mathrm{sim}}(k)
  \right]^2 .
  \label{eq:msd-tobs}
\end{equation}
For each simulated trajectory $m$, the reference statistic is computed against
the leave-one-out simulated mean,
\begin{equation}
  T_m^{\mathrm{MSD}}
  =
  \sum_{k \in \mathcal{K}}
  w_k
  \left[
    \mathrm{MSD}_m(k)
    -
    \overline{\mathrm{MSD}}_{-m}(k)
  \right]^2 ,
  \label{eq:msd-tsim}
\end{equation}
where
\[
  \overline{\mathrm{MSD}}_{-m}(k)
  =
  \frac{1}{M-1}\sum_{\ell \ne m}\mathrm{MSD}_{\ell}(k).
\]
This discrepancy is sensitive to both the magnitude and the shape of the MSD
curve. It can detect trajectories that diffuse too rapidly, remain too
spatially constrained, or exhibit non-smooth displacement dynamics that are not
reproduced by the fitted HMM-SSF. In multi-state models, MSD is particularly
useful because it depends jointly on state-specific movement kernels and on the
temporal arrangement of states. A model may match average step lengths within
states but still fail to reproduce MSD if residence times or switching
patterns are wrong.

\paragraph{Sinuosity discrepancy.}
Path sinuosity targets the geometry of the realized path rather than its
spatial footprint or displacement scaling. We use the straightness index
\citep{benhamou2004},
\begin{equation*}
  \mathrm{SI}_r
  =
  \frac{\|X_T^{(r)}-X_1^{(r)}\|}
  {\sum_{t=1}^{T-1}\|X_{t+1}^{(r)}-X_t^{(r)}\|},
  \qquad 0 \le \mathrm{SI}_r \le 1.
\end{equation*}
Low values correspond to more tortuous paths, whereas values close to one
indicate comparatively straight movement. Let
\[
  \overline{\mathrm{SI}}_{\mathrm{sim}}
  =
  \frac{1}{M}\sum_{m=1}^{M}\mathrm{SI}_m
\]
be the simulated mean straightness index. The observed discrepancy is
\begin{equation*}
  T_{\mathrm{obs}}^{\mathrm{SI}}
  =
  \left|
    \mathrm{SI}_0
    -
    \overline{\mathrm{SI}}_{\mathrm{sim}}
  \right|,
\end{equation*}
and the simulated reference discrepancies are
\begin{equation*}
  T_m^{\mathrm{SI}}
  =
  \left|
    \mathrm{SI}_m
    -
    \overline{\mathrm{SI}}_{-m}
  \right|,
  \qquad
  \overline{\mathrm{SI}}_{-m}
  =
  \frac{1}{M-1}\sum_{\ell \ne m}\mathrm{SI}_{\ell}.
\end{equation*}
The absolute discrepancy is used because both directions are biologically
meaningful: simulated paths can be too straight, suggesting excessive
directional persistence or too much time in travelling-like regimes, or too
tortuous, suggesting excessive turning or too much time in localized regimes.
In HMM-SSFs, sinuosity is therefore a useful bridge between movement kernels
and state occupancy: it can flag cases where the marginal spatial distribution
looks acceptable but the behavioural geometry of the path is unrealistic.

\paragraph{Barrier-crossing discrepancy.}
When a known barrier geometry is available, the HMM-SSF workflow inherits the
barrier-crossing pillar from the single-state SSF framework. The diagnostic
compares the observed number of intersections between movement steps and a
specified linear feature with the corresponding simulated distribution. Let
$B$ denote the barrier geometry and let
\begin{equation*}
  C(X_{1:T}^{(r)};B)
  =
  \sum_{t=1}^{T-1}
  \ind\{[X_t^{(r)},X_{t+1}^{(r)}]\cap B \ne \emptyset\}
\end{equation*}
be the number of step segments crossing the barrier. When the scientific
concern is excess permeability, the one-sided Monte Carlo p-value is
\begin{equation}
  p_{\mathrm{barrier}}
  =
  \frac{
    1+
    \sum_{m=1}^{M}
    \ind\{
      C(X_{1:T}^{(m)};B)
      \le
      C(X_{1:T}^{(0)};B)
    \}
  }{M+1}.
  \label{eq:barrier-pvalue}
\end{equation}
Small values indicate that the observed trajectory crosses the barrier less
often than expected under the fitted model, suggesting that the model
overestimates permeability. If the ecological concern is instead that the
model is too restrictive, the inequality in Eq.~\eqref{eq:barrier-pvalue} can
be reversed. In a latent-state model, barrier mismatch can arise from the
state-specific selection kernels, from unrealistic residence near the barrier,
or from the state process assigning barrier-crossing opportunities to the wrong
regimes. This diagnostic is therefore interpreted jointly with state
occupancy, residence time, and state-conditioned movement summaries when
state-specific barrier response is scientifically relevant. As in the
single-state framework, the barrier must be specified a priori and the
diagnostic can have low power when crossing events are rare.

\subsection{State-process diagnostics}

State-process diagnostics require an observed or decoded state path and a
corresponding simulated state path. In empirical applications, observed states
are typically decoded, for example using the Viterbi path, rather than directly
observed. The diagnostics should therefore be interpreted as conditional on a
state-assignment procedure, or evaluated across several state-simulation
strategies to assess sensitivity to latent-state uncertainty. Let
$A_{1:T}^{(0)}$ denote the decoded or supplied state path associated with the
observed trajectory, and let $A_{1:T}^{(m)}$, $m=1,\ldots,M$, denote simulated
state paths generated under the fitted HMM-SSF. The following diagnostics ask
whether $A_{1:T}^{(0)}$ is typical of the state-process distribution induced by
the fitted model, either marginally or after conditioning on the chosen
state-assignment strategy.

\paragraph{State occupancy.}
State occupancy measures the time budget allocated to each behavioural regime.
For a state sequence $a_{1:T}$, define the occupancy proportion of state $s$ as
\begin{equation*}
  \pi_s(a)
  =
  \frac{1}{T}\sum_{t=1}^{T}\ind\{a_t=s\}.
\end{equation*}
For each state, the observed occupancy discrepancy is the absolute difference
between the decoded observed occupancy and the mean simulated occupancy:
\begin{equation*}
  T_{\mathrm{obs},s}^{\mathrm{occ}}
  =
  \left|
    \pi_s(A^{(0)})
    -
    \frac{1}{M}\sum_{m=1}^{M}\pi_s(A^{(m)})
  \right|.
\end{equation*}
The simulated reference discrepancies are obtained by leave-one-out comparison:
\begin{equation*}
  T_{m,s}^{\mathrm{occ}}
  =
  \left|
    \pi_s(A^{(m)})
    -
    \frac{1}{M-1}\sum_{\ell \ne m}\pi_s(A^{(\ell)})
  \right|.
\end{equation*}
A global occupancy summary is also useful:
\begin{equation*}
  T_{\mathrm{obs}}^{\mathrm{occ,global}}
  =
  \frac{1}{K}
  \sum_{s=1}^{K}
  \left|
    \pi_s(A^{(0)})
    -
    \frac{1}{M}\sum_{m=1}^{M}\pi_s(A^{(m)})
  \right|.
\end{equation*}
For two-state models, per-state absolute discrepancies can be identical because
the two occupancy proportions are complementary. In such cases, the global
occupancy statistic provides the model-level interpretation, while per-state
rows identify which biological regimes are over- or under-represented.

\paragraph{Residence time.}
Residence times are run lengths of consecutive visits to the same state. Let
\[
  \mathcal{R}_s(a)
  =
  \{r_{s,1}(a),\ldots,r_{s,n_s(a)}(a)\}
\]
be the empirical distribution of run lengths in state $s$ for sequence $a$.
The simulation-based residence-time diagnostic compares
$\mathcal{R}_s(A^{(0)})$ with the corresponding simulated distributions using a
one-dimensional Wasserstein distance:
\begin{equation}
  T_{\mathrm{obs},s}^{\mathrm{res}}
  =
  \frac{1}{M}
  \sum_{m=1}^{M}
  W_1\{\mathcal{R}_s(A^{(0)}),\mathcal{R}_s(A^{(m)})\}.
  \label{eq:res-tobs}
\end{equation}
The simulated reference statistic is
\begin{equation}
  T_{m,s}^{\mathrm{res}}
  =
  \frac{1}{M-1}
  \sum_{\ell \ne m}
  W_1\{\mathcal{R}_s(A^{(m)}),\mathcal{R}_s(A^{(\ell)})\}.
  \label{eq:res-tsim}
\end{equation}
A small p-value indicates that residence times in the decoded observed state
path are unusual relative to residence times generated by the fitted transition
model. This is a stronger statement than an occupancy discrepancy: a model can
reproduce the overall time budget in each state while producing bouts that are
too short, too long, or too variable.

\paragraph{Geometric residence-time check.}
In a homogeneous HMM, the residence time in state $s$ has a geometric
distribution with leaving probability $1-\Gamma_{ss}$:
\begin{equation}
  \Pr(R_s=r)
  =
  \Gamma_{ss}^{r-1}(1-\Gamma_{ss}),
  \qquad r=1,2,\ldots .
  \label{eq:geometric-dwell}
\end{equation}
This property follows directly from the Markov assumption: once the process is
in state $s$, the probability of leaving at the next time step is constant and
does not depend on how long the process has already remained in that state.
The geometric residence-time diagnostic targets this assumption directly. For
each state, bootstrap samples
$G_s^{(1)},\ldots,G_s^{(B)}$ are generated from the geometric distribution in
Eq.~\eqref{eq:geometric-dwell}, with sample size matching the number of
observed residence bouts in state $s$. The observed discrepancy is
\begin{equation*}
  T_{\mathrm{obs},s}^{\mathrm{geo}}
  =
  \frac{1}{B}\sum_{b=1}^{B}
  W_1\{\mathcal{R}_s(A^{(0)}),G_s^{(b)}\},
\end{equation*}
and the parametric-bootstrap reference uses leave-one-out distances among the
geometric samples:
\begin{equation*}
  T_{b,s}^{\mathrm{geo}}
  =
  \frac{1}{B-1}\sum_{c \ne b}
  W_1(G_s^{(b)},G_s^{(c)}).
\end{equation*}
A small p-value indicates that the decoded dwell-time distribution is unusual
under the geometric distribution implied by the fitted homogeneous HMM. This is
evidence of a residence-time mismatch, not definitive proof that a semi-Markov
model is the correct alternative. The diagnostic is defined only for
homogeneous transition matrices; for time-varying transition probabilities,
residence-time adequacy should be assessed using model-generated state paths as
in Eqs.~\eqref{eq:res-tobs}--\eqref{eq:res-tsim}.

\paragraph{Switching rate and transition counts.}
The switching rate summarizes the overall frequency of changes between latent
regimes:
\begin{equation*}
  \rho(a)
  =
  \frac{1}{T-1}
  \sum_{t=1}^{T-1}\ind\{a_{t+1}\ne a_t\}.
\end{equation*}
The observed switching discrepancy is
\begin{equation*}
  T_{\mathrm{obs}}^{\mathrm{switch}}
  =
  \left|
    \rho(A^{(0)})
    -
    \frac{1}{M}\sum_{m=1}^{M}\rho(A^{(m)})
  \right|,
\end{equation*}
with leave-one-out simulated discrepancies defined analogously. This diagnostic
detects whether the fitted model switches too often or not often enough on
average, but it does not identify which transitions are responsible.

To retain directional information, let
\[
  N_{ij}(a)
  =
  \sum_{t=1}^{T-1}\ind\{a_t=i, a_{t+1}=j\}
\]
be the empirical transition-count matrix, and define its normalized version
\[
  P_{ij}(a)
  =
  \frac{N_{ij}(a)}{\sum_{u=1}^{K}\sum_{v=1}^{K}N_{uv}(a)}.
\]
Writing $\mathrm{vec}\{P(a)\}$ for the vectorized normalized transition matrix,
the transition-count discrepancy is
\begin{equation*}
  T_{\mathrm{obs}}^{\mathrm{trans}}
  =
  \left\|
    \mathrm{vec}\{P(A^{(0)})\}
    -
    \frac{1}{M}
    \sum_{m=1}^{M}
    \mathrm{vec}\{P(A^{(m)})\}
  \right\|_2,
\end{equation*}
with simulated reference discrepancies computed by leave-one-out replacement
of the simulated mean. Transition-count diagnostics can agree with the fitted
HMM even when residence-time shape is misspecified, because a Markov chain can
reproduce average transition frequencies without reproducing non-geometric
dwell-time structure.

\paragraph{State-conditioned movement summaries.}
State-conditioned diagnostics compare movement summaries within each decoded or
simulated state. We consider state-conditioned mean squared displacement,
step-length distributions, and turning-angle summaries. These diagnostics
target a different failure mode from occupancy or residence time: they ask
whether the state-specific movement kernels reproduce the movement phenotype
associated with each regime, conditional on the chosen state assignment.

For step lengths, steps are assigned to the state at the arrival location. Let
$L_s^{(r)}$ denote the empirical distribution of step lengths assigned to state
$s$ in trajectory $r$. The state-specific step-length discrepancy is
\begin{equation*}
  T_{\mathrm{obs},s}^{\mathrm{step}}
  =
  \frac{1}{M}
  \sum_{m=1}^{M}
  W_1(L_s^{(0)},L_s^{(m)}),
\end{equation*}
with leave-one-out Wasserstein distances among simulated state-specific
step-length distributions used as the reference. A small p-value indicates
that the movement kernel for state $s$ does not reproduce the observed
state-specific step-length distribution.

For turning angles, we summarize directional concentration in state $s$ using
the mean resultant length
\begin{equation*}
  R_s
  =
  \sqrt{
    \left[
      \frac{1}{n_s}\sum_{q=1}^{n_s}\cos(\alpha_{s,q})
    \right]^2
    +
    \left[
      \frac{1}{n_s}\sum_{q=1}^{n_s}\sin(\alpha_{s,q})
    \right]^2
  },
\end{equation*}
where $\alpha_{s,q}$ are the turning angles assigned to state $s$. Values near
zero indicate dispersed turning angles, while values near one indicate strong
directional persistence. The diagnostic compares the observed $R_s$ with the
mean simulated $R_s$ and uses leave-one-out simulated discrepancies as the
reference.

For within-state MSD, locations assigned to state $s$ are extracted in temporal
order and treated as a state-specific subsequence. An MSD curve is then
computed on this ordered subsequence and compared with the mean simulated
state-specific MSD curve using the same integrated squared discrepancy as in
Eqs.~\eqref{eq:msd-tobs}--\eqref{eq:msd-tsim}. This diagnostic is descriptive,
because the state-specific subsequence need not have evenly spaced dwell-time
structure after extraction, but it is useful for identifying regimes whose
within-state displacement dynamics are not reproduced by the fitted movement
kernel.

\section{Modelling and prediction with HMM-SSFs}

The generative representation in Eq.~\eqref{eq:joint-generative} makes an
HMM-SSF useful not only for retrospective behavioural interpretation, but also
for prediction. Prediction is not a single task, however. Depending on the
scientific question, the target may be the next step, the next behavioural
state, a full future trajectory, an emergent utilization distribution, a
barrier-crossing frequency, or a counterfactual response to a changed
landscape. These targets use the same fitted model but place different demands
on the transition model, the state-specific movement kernels, and the
state-specific selection functions.

\begin{table}[ht]
\centering
\caption{Prediction targets induced by a fitted HMM-SSF.}
\label{tab:prediction-targets}
\begin{tabular}{p{0.22\linewidth}p{0.34\linewidth}p{0.34\linewidth}}
\toprule
Prediction target & Model component emphasized & Typical output \\
\midrule
One-step movement & Filtered state probabilities and state-specific SSF kernels & Conditional surface or candidate-step probabilities for $X_{t+1}$ \\
State forecasting & Transition matrix or transition model $\Gamma(z_t)$ & Predicted state probabilities, occupancy, switching probability \\
Trajectory forecasting & Full latent-state and movement process & Simulated paths over a specified horizon \\
Emergent space use & Iterated movement process and landscape covariates & Predicted utilization distribution or occupancy surface \\
Connectivity or barrier response & State-specific selection and movement near known barriers & Crossing counts, first-passage summaries, side-specific residence \\
Scenario prediction & Modified covariates or transition drivers & Counterfactual trajectories and derived summaries \\
\bottomrule
\end{tabular}
\end{table}

\subsection{From fitted model to predictive distribution}

At short horizons, prediction can be expressed analytically as mixtures of
state-specific kernels. Eq.~\eqref{eq:one-step-predictive} shows that the
one-step predictive distribution combines three elements: the current state
uncertainty $\alpha_t$, the transition probabilities $\Gamma(z_t)$, and the
state-specific SSF kernels $f_j$. This decomposition is useful for
interpretation. A predicted move can be uncertain because the next behavioural
state is uncertain, because the state-specific movement kernel is diffuse, or
because habitat selection assigns similar weights to many available endpoints.

At longer horizons, exact integration is generally impractical because future
locations determine future movement covariates, available steps, and
environmental values. Forward simulation is therefore the natural predictive
tool. A single predictive replicate is generated by repeating the following
steps: sample the next state from $\Gamma(z_t)$, sample the next location from
the corresponding state-specific SSF, update the movement history, update any
time-varying covariates, and continue until the prediction horizon is reached.
An ensemble of such replicates approximates the predictive distribution of any
trajectory-level summary.

\subsection{Latent-state simulation as predictive conditioning}

The hidden component of an HMM-SSF is a discrete stochastic process. Simulating
it means generating a state sequence
\[
  S_{1:H}^{\ast} = (S_1^{\ast},\ldots,S_H^{\ast})
\]
over a prediction or validation horizon $H$, then using this sequence to choose
which state-specific SSF kernel generates each movement step. For the full
generative model, the simulated state path is drawn from the fitted initial
distribution and transition model:
\begin{align*}
  S_1^{\ast} &\sim \mathrm{Categorical}(\delta), \\
  S_{h+1}^{\ast} \given S_h^{\ast}=i, z_h^{\ast}
  &\sim
  \mathrm{Categorical}\{\Gamma_{i1}(z_h^{\ast}),\ldots,\Gamma_{iK}(z_h^{\ast})\},
  \qquad h=1,\ldots,H-1 .
\end{align*}
Equivalently, if $U_h \sim \mathrm{Uniform}(0,1)$, then
\begin{equation*}
  S_{h+1}^{\ast}
  =
  \min\left\{
    j:
    \sum_{\ell=1}^{j} \Gamma_{i\ell}(z_h^{\ast}) \ge U_h
  \right\},
  \qquad i=S_h^{\ast}.
\end{equation*}
This inverse-transform expression is useful because it makes clear that the
next hidden state is sampled from the row of the transition matrix
corresponding to the current hidden state.

In a homogeneous HMM, $\Gamma(z_h^{\ast})=\Gamma$ for all $h$, so the hidden
state simulation is an ordinary finite-state Markov chain. In a time-varying
HMM, the transition matrix is recomputed at each step from the relevant
covariates:
\[
  \Gamma(z_h^{\ast})
  =
  \left[
    \Gamma_{ij}(z_h^{\ast})
  \right]_{i,j=1}^{K}.
\]
These covariates may be external, such as time of day or season, or they may be
functions of the simulated trajectory itself. When transition covariates depend
on the simulated trajectory, the hidden and spatial components must be updated
jointly: draw $S_{h+1}^{\ast}$, draw $X_{h+1}^{\ast}$ from the corresponding
state-specific SSF, update the covariates, then proceed to the next step.

Conditional on the simulated hidden path, the simulated spatial path satisfies
\begin{equation*}
  X_{h+1}^{\ast}
  \given X_{h}^{\ast}, X_{h-1}^{\ast}, S_{h+1}^{\ast}=s
  \sim
  f_s(\cdot \given X_h^{\ast}, X_{h-1}^{\ast}, e_h^{\ast}).
\end{equation*}
Thus the full predictive simulation can be written as the product
\begin{equation*}
  p_{\theta}(s_{1:H}^{\ast},x_{1:H}^{\ast})
  =
  p_{\theta}(s_{1:H}^{\ast})
  p_{\theta}(x_{1:H}^{\ast}\given s_{1:H}^{\ast}),
\end{equation*}
where the first term is the hidden Markov chain and the second term is the
sequence of state-conditioned SSF draws. This decomposition is the basis for
the different simulation strategies below.

We distinguish four simulation strategies. They are not competing fitted
models; they are four ways of conditioning the same fitted HMM-SSF before
simulating the spatial trajectory. The difference among them is entirely in
how the latent path $S_{1:H}^{\ast}$ is generated.

\subsubsection*{Markov generation}

Markov generation is the full generative use of the fitted HMM-SSF. The latent
state sequence is simulated from the fitted initial distribution and fitted
transition model:
\begin{align*}
  S_1^{\ast}
  &\sim
  \mathrm{Categorical}(\hat{\delta}), \\
  S_{h+1}^{\ast}\given S_h^{\ast}=i
  &\sim
  \mathrm{Categorical}
  \{\hat{\Gamma}_{i1}(z_h^{\ast}),\ldots,\hat{\Gamma}_{iK}(z_h^{\ast})\}.
\end{align*}
The spatial path is then generated conditionally on this sampled state path:
\begin{equation*}
  X_{h+1}^{\ast}
  \given X_h^{\ast},X_{h-1}^{\ast},S_{h+1}^{\ast}=s
  \sim
  f_s(\cdot \given X_h^{\ast},X_{h-1}^{\ast},e_h^{\ast};\hat{\theta}).
\end{equation*}
Equivalently, this strategy samples from
\[
  p_{\hat{\theta}}(s_{1:H}^{\ast},x_{1:H}^{\ast})
  =
  p_{\hat{\theta}}(s_{1:H}^{\ast})
  p_{\hat{\theta}}(x_{1:H}^{\ast}\given s_{1:H}^{\ast}).
\]
It is the appropriate strategy for genuine forward prediction, simulation of
space use, and scenario analysis because it propagates both latent-state
uncertainty and movement stochasticity.

\subsubsection*{Viterbi-conditioned generation}

Viterbi-conditioned generation fixes the latent path to the most likely decoded
state sequence for the observed trajectory. If
\begin{equation*}
  \hat{s}_{1:T}^{V}
  =
  \arg\max_{s_{1:T}}
  p_{\hat{\theta}}(s_{1:T}\given x_{1:T}),
\end{equation*}
then this strategy sets
\begin{equation*}
  S_{1:T}^{\ast}=\hat{s}_{1:T}^{V}
\end{equation*}
and simulates only the spatial component
\begin{equation*}
  X_{1:T}^{\ast}
  \sim
  p_{\hat{\theta}}(X_{1:T}\given S_{1:T}=\hat{s}_{1:T}^{V}).
\end{equation*}
This is not a full predictive distribution for future movement, because the
latent state sequence is treated as known. Its purpose is diagnostic: it asks
whether the state-specific movement and selection kernels can reproduce the
trajectory when the regime sequence is held fixed. If the Viterbi-conditioned
strategy is compatible with the observed trajectory but Markov generation is
not, the discrepancy is more plausibly associated with the transition model,
state occupancy, residence times, or switching structure.

\subsubsection*{Posterior state-path generation}

Posterior state-path generation samples the latent path from the smoothing
distribution implied by the fitted model and observed trajectory:
\begin{equation*}
  S_{1:T}^{\ast}
  \sim
  p_{\hat{\theta}}(S_{1:T}\given X_{1:T}=x_{1:T}).
\end{equation*}
In practice, this can be implemented by forward-filtering backward-sampling.
If $\alpha_h(i)=\Pr_{\hat{\theta}}(S_h=i\given X_{1:h}=x_{1:h})$ denotes the
filtered probability at time $h$, then a backward draw satisfies
\begin{align*}
  S_T^{\ast}
  &\sim
  \Pr_{\hat{\theta}}(S_T\given X_{1:T}=x_{1:T}), \\
  \Pr(S_h^{\ast}=i\given S_{h+1}^{\ast}=j,X_{1:T}=x_{1:T})
  &\propto
  \alpha_h(i)\hat{\Gamma}_{ij}(z_h),
  \qquad h=T-1,\ldots,1 .
\end{align*}
After a posterior state path is sampled, the spatial trajectory is simulated
from
\[
  X_{1:T}^{\ast}
  \sim
  p_{\hat{\theta}}(X_{1:T}\given S_{1:T}^{\ast}).
\]
This strategy incorporates uncertainty in state assignment while remaining
conditioned on the information in the observed trajectory. It is useful when
the scientific question is whether diagnostic conclusions depend on treating a
single decoded state path as fixed.

\subsubsection*{Viterbi-tube generation}

Viterbi-tube generation is a sensitivity strategy centred on the Viterbi path.
Let $d(s_{1:T},\hat{s}_{1:T}^{V})$ be a distance between state sequences, for
example a Hamming distance or another application-specific discrepancy. For a
tolerance $\epsilon$, define the tube around the Viterbi path as
\begin{equation*}
  \mathcal{T}_{\epsilon}(\hat{s}^{V})
  =
  \{s_{1:T}: d(s_{1:T},\hat{s}_{1:T}^{V})\le \epsilon\}.
\end{equation*}
The idealized Viterbi-tube strategy samples from the posterior distribution
restricted to this tube:
\begin{equation*}
  \Pr_{\hat{\theta}}(S_{1:T}^{\ast}=s_{1:T}
  \given X_{1:T}=x_{1:T}, S_{1:T}\in\mathcal{T}_{\epsilon})
  =
  \frac{
    p_{\hat{\theta}}(s_{1:T}\given x_{1:T})
    \ind\{s_{1:T}\in\mathcal{T}_{\epsilon}(\hat{s}^{V})\}
  }{
    \sum_{\tilde{s}_{1:T}\in\mathcal{T}_{\epsilon}(\hat{s}^{V})}
    p_{\hat{\theta}}(\tilde{s}_{1:T}\given x_{1:T})
  }.
\end{equation*}
The spatial trajectory is then simulated conditionally on the sampled path.
This strategy is not intended as a new ecological model. It asks whether the
diagnostic conclusions depend on one exact Viterbi sequence or remain stable
for nearby high-support regime sequences.

\subsection{Prediction after model checking}

The diagnostics in this paper are intended to make prediction more explicit.
For example, if utilization-distribution diagnostics are compatible but MSD
diagnostics fail, the model may still be useful for long-term spatial footprint
prediction while being unreliable for transient displacement or encounter-rate
prediction. If state occupancy is compatible but residence-time diagnostics
fail, predicted time budgets may be reasonable while predicted bout durations
or switching sequences remain unreliable. If state-conditioned step-length
diagnostics fail, forecasts that depend on travel distance, connectivity, or
barrier-crossing opportunities should be treated cautiously even if global
space-use summaries look acceptable.

This perspective separates model fitting from model use. A fitted HMM-SSF can
be statistically estimable and biologically interpretable without being
adequate for every predictive purpose. Generative validation identifies which
features of the predictive distribution are supported by the fitted model and
which features require model refinement, additional states, time-varying
transition structure, memory, or richer state-specific movement kernels.

\section{Simulation study}

This section defines the proposed simulation programme. No simulation results
are reported at this stage. The objective is to specify a design that can be
discussed before implementation, so that the final study tests the
methodological claims of the paper rather than only illustrating the software.

The simulation study has four goals. First, it should assess calibration under
well-specified HMM-SSFs: when the fitted model matches the data-generating
process, the observed trajectory should not be systematically extreme relative
to the fitted generative distribution. Second, it should evaluate diagnostic
sensitivity under controlled misspecification, with each scenario targeting a
distinct model component. Third, it should determine whether trajectory-level
and state-process diagnostics identify different failure modes. Fourth, it
should clarify how conclusions depend on the way latent state uncertainty is
handled during simulation.

\paragraph{General protocol.}
All scenarios will use the same basic workflow. A known data-generating process
will be used to simulate a movement trajectory, the candidate model or models
will be fitted to that trajectory, and the fitted model will be used to
generate replicated trajectories and state sequences. The diagnostic workflow
will then compare the observed simulated trajectory with the model-generated
replicates. For each scenario, the planned outputs are Monte Carlo p-values,
observed discrepancies, simulation-reference variability, diagnostic sharpness,
and rejection frequencies across independent replicate data sets.

The default synthetic landscape will contain one smooth habitat covariate and,
for the barrier scenarios, one linear barrier. The default behavioural
interpretation will use an encamped state with short, tortuous steps and a
travelling state with longer, more persistent movement. Step lengths will be
generated from state-specific positive distributions, turning angles from
state-specific circular distributions, and endpoint weights from
state-specific selection coefficients. Unless a scenario explicitly changes
one component, the fitted model will contain the same covariates and the same
number of states as the data-generating process.

\begin{table}[ht]
\centering
\caption{Planned simulation scenarios for evaluating HMM-SSF generative diagnostics.}
\label{tab:simulation-plan}
\renewcommand{\arraystretch}{1.15}
\setlength{\tabcolsep}{3pt}
\footnotesize
\begin{tabular}{p{0.08\linewidth}p{0.22\linewidth}p{0.27\linewidth}p{0.31\linewidth}}
\toprule
ID & Purpose & Fitted models & Primary diagnostic target \\
\midrule
S0 & Calibration under a correct two-state HMM-SSF. &
Correct two-state HMM-SSF. &
Approximate calibration of all diagnostic p-values. \\

S1 & Loss induced by collapsing two behavioural regimes into one. &
One-state SSF and correct two-state HMM-SSF. &
Trajectory-level MSD and sinuosity, with UD used as a contrasting spatial diagnostic. \\

S2 & Incorrect number of latent states. &
One-, two-, three-, and four-state HMM-SSFs fitted to data generated with three states. &
Occupancy, state-conditioned movement, and trajectory geometry. \\

S3 & Missing time variation in the transition matrix. &
Fixed-TPM and time-varying-TPM HMM-SSFs. &
Occupancy, switching rate, residence time, and MSD. \\

S4 & Misspecified state-specific movement kernel. &
Simplified movement kernel compared with the correct movement kernel. &
State-specific step length, turning angle, within-state MSD, and global MSD. \\

S5 & Misspecified habitat-selection structure. &
No habitat effect, shared habitat effect, and correct state-specific habitat effect. &
UD discrepancy and state-conditioned movement or space-use summaries. \\

S6 & Misspecified barrier response. &
No barrier response, common barrier response, and correct state-specific barrier response. &
Barrier crossing and state-specific crossing opportunities. \\

S7 & Non-geometric residence times. &
HMM-SSF fitted to trajectories generated from an HSMM-like state process. &
Residence-time and geometric-dwell diagnostics, with occupancy as a negative control. \\

S8 & Sensitivity to state uncertainty and weak state separation. &
Correct HMM-SSF analyzed with Markov, Viterbi-conditioned, posterior path, and Viterbi-tube simulations. &
Robustness of state-process diagnostics to the state-assignment rule. \\
\bottomrule
\end{tabular}
\end{table}

\FloatBarrier

\paragraph{Scenario S0: well-specified two-state HMM-SSF.}
The data-generating process will be a two-state HMM-SSF with one encamped state
and one travelling state. The two states will differ in step-length scale,
turning-angle concentration, habitat-selection strength, and persistence. The
fitted model will use the correct number of states, the correct habitat
covariate, and the correct transition structure. This scenario is the
calibration check for the entire workflow. The expected result is that
trajectory-level and state-process p-values are not systematically small,
apart from finite-sample and state-decoding effects.

\paragraph{Scenario S1: single-state approximation to a two-state process.}
The data-generating process will be the same as in S0, but the fitted
comparison will include a one-state SSF and the correctly specified two-state
HMM-SSF. This scenario tests the main motivation for the HMM extension of the
single-state framework. The one-state model may reproduce parts of the spatial
footprint, especially when the two states use similar areas, but it should
struggle to reproduce displacement scaling and path geometry. The primary
diagnostics are MSD and sinuosity; UD is included to show that spatial
compatibility alone does not imply dynamic adequacy.

\paragraph{Scenario S2: incorrect number of states.}
The data-generating process will contain three states: an encamped state, an
intermediate movement state, and a travelling state. Candidate models with one,
two, three, and four states will be fitted. The purpose is to distinguish
underfitting from overfitting. Models with too few states should show
state-conditioned movement discrepancies, distorted occupancy, and possible
MSD or sinuosity failures. Models with too many states may generate redundant
or poorly occupied regimes, which should appear in occupancy, transition, and
state-conditioned diagnostics.

\paragraph{Scenario S3: missing time-varying transition structure.}
The data-generating process will use a two-state HMM-SSF in which the
transition matrix depends on a cyclic temporal covariate, such as time of day.
The fitted comparison will include a fixed-TPM model and the correct
time-varying-TPM model. The goal is to test whether the diagnostics detect
temporal organization in the state process that is not visible from the
state-specific movement kernels alone. The primary diagnostics are occupancy
through time, switching rate, residence time, transition counts, and MSD.

\paragraph{Scenario S4: misspecified movement kernel.}
The data-generating process will keep the correct state-transition and
selection structure, but one state will have a movement distribution that is
not available to the fitted simplified model. For example, the travelling
state may have heavier-tailed step lengths or stronger directional persistence
than the fitted kernel allows. The fitted comparison will include the
simplified movement kernel and, if available, the correctly specified movement
kernel. This scenario targets state-specific step-length discrepancies,
turning-angle concentration, within-state MSD, and global MSD.

\paragraph{Scenario S5: misspecified habitat selection.}
The data-generating process will include state-specific habitat selection. For
example, the encamped state may strongly select high values of the habitat
covariate, while the travelling state may respond weakly or in the opposite
direction. Three fitted models will be compared: a model with no habitat
effect, a model with a shared habitat coefficient across states, and the
correct state-specific model. This scenario tests whether the UD diagnostic
and state-conditioned summaries can detect selection misspecification even
when movement kernels and transition probabilities are otherwise adequate.

\paragraph{Scenario S6: misspecified barrier response.}
The data-generating process will include a linear barrier that affects
movement differently by state. For example, the encamped state may rarely
cross the barrier, while the travelling state may cross it with lower but
nonzero resistance. Candidate fitted models will omit the barrier, include a
common barrier effect, or include the correct state-specific barrier response.
The primary diagnostic is the barrier-crossing statistic, interpreted jointly
with state occupancy near the barrier and state-conditioned crossing
opportunities.

\paragraph{Scenario S7: residence-time memory.}
The data-generating process will use an HSMM-like state sequence with
non-geometric residence times, while keeping HMM-like state-specific movement
and selection kernels. The fitted model will be an HMM-SSF, so the fitted
state process will impose geometric residence times. This scenario is designed
to answer whether the residence-time diagnostics can detect memory in the
latent process. The expected signature is that residence-time and
geometric-dwell diagnostics reject, while occupancy and average switching rate
may remain compatible.

\paragraph{Scenario S8: state uncertainty and weak state separation.}
The data-generating process will be a correctly specified two-state HMM-SSF
with weak separation between states. Two variants will be used: one where
states differ mainly in movement and one where states differ mainly in habitat
selection. The fitted model will be correct, but diagnostics will be computed
using true simulated states, the Viterbi path, posterior state-path samples,
and Viterbi-tube simulations. This scenario tests whether state-process
conclusions are robust to state uncertainty or depend on a single decoded path.

\paragraph{Reporting structure.}
The final simulation section should report calibration and power separately.
Calibration will be assessed from S0 using the distribution of p-values across
replicate data sets. Diagnostic sensitivity will be summarized from S1--S8 as
rejection frequencies by scenario and diagnostic family. Each scenario should
also include one representative diagnostic plot, chosen to show the mechanism
of failure rather than only the numerical p-value. Before implementation, the
remaining fixed choices are the number of replicate data sets, the trajectory
lengths, the number of available steps, and the Monte Carlo simulation count
used inside each diagnostic test.

\FloatBarrier

\section{Zebra model-comparison example}

The empirical demonstration uses the zebra HMM-SSF example of
\citet{klappstein2023}. Four candidate models are compared using Markov
generation for all fitted models:

\begin{table}[ht]
\centering
\caption{Candidate models in the zebra comparison.}
\label{tab:zebra-models}
\begin{tabular}{llll}
\toprule
Label & States & Transition model & Interpretation \\
\midrule
M1 & 1 & Degenerate & Single-state SSF baseline \\
M2 & 2 & Constant TPM & Two states, no time-of-day effect \\
M3 & 2 & Time-varying TPM & Two states, time-of-day effect \\
M4 & 3 & Constant TPM & Three-state fixed-transition alternative \\
\bottomrule
\end{tabular}
\end{table}

The SSF formula in the zebra application is
\[
  \texttt{\char`~ step + log(step) + cos(angle) + veg},
\]
which includes movement terms and a vegetation covariate. The time-varying
transition model uses cyclic time-of-day terms,
\[
  \texttt{\char`~ cos(2*pi*tod/24) + sin(2*pi*tod/24)}.
\]
The comparison uses 99 Markov simulations per candidate model, giving a Monte
Carlo p-value resolution of 0.01.

\begin{table}[ht]
\centering
\caption{Trajectory-level zebra diagnostics from the current empirical analysis. Values are observed discrepancies with Monte Carlo p-values in parentheses.}
\label{tab:zebra-global}
\begin{tabular}{llll}
\toprule
Model & UD Wasserstein & MSD & Sinuosity \\
\midrule
M1 one-state SSF & 19.04 (0.61) & $1.25\times 10^6$ (0.02) & 0.175 (0.01) \\
M2 two-state fixed TPM & 12.99 (0.30) & $3.48\times 10^4$ (0.01) & 0.108 (0.05) \\
M3 two-state time TPM & 12.81 (0.25) & $3.66\times 10^4$ (0.01) & 0.107 (0.06) \\
M4 three-state fixed TPM & 12.91 (0.24) & $2.95\times 10^4$ (0.01) & 0.099 (0.13) \\
\bottomrule
\end{tabular}
\end{table}

The UD discrepancy is not extreme for any candidate model in this run. MSD is
extreme for every candidate model, including the multi-state alternatives. The
one-state model fails strongly on sinuosity, whereas the multi-state models are
less discrepant, with the three-state fixed model giving the largest sinuosity
p-value in this analysis.

\begin{figure}[ht]
\centering
\maybeincludegraphics[width=0.82\linewidth]{../docs/articles/zebra-model-comparison_files/figure-html/global-pvalue-heatmap-1.png}
\caption{Monte Carlo p-values for the trajectory-level zebra diagnostics.}
\label{fig:zebra-pvalues}
\end{figure}

State-level diagnostics are interpreted conditionally on the decoded observed
state path. In two-state fits, per-state occupancy discrepancies can be
identical because the occupancy proportions sum to one. The global occupancy
discrepancy is therefore preferred for model-level reporting, with per-state
rows used to identify which regime contributes to the mismatch.

\bibliographystyle{apalike}
\bibliography{references}

\end{document}
